\documentclass{article}

\usepackage{tikz}

\makeatletter
% 1. 主命令：负责接收可选参数 [] 和圆括号内容 ()
\newcommand{\segment}[2][]{%
  \message{#2}
  \segment@impl[#1]#2%
}

\def\segment@impl[#1](#2){
  \pgfutil@in@{,}{#2}
  \ifpgfutil@in@%
    % 如果括号内包含逗号，调用处理双点的宏
    \segmentAB[#1](#2)%
  \else%
    % 如果不包含逗号，调用处理单点/标签的宏
    \segmentL[#1](#2)%
  \fi%
}

% 2. 辅助宏 A：处理 \segment[options](A,B)
% 注意这里的参数定界符 (#2,#3)
\def\segmentAB[#1]#2(#3,#4){%
  \draw[#1] (#3) -- (#4);
}

% 3. 辅助宏 B：处理 \segment[options](L)
\def\segmentL[#1]#2(#3){%
  \fill[#1] (0,0) -- (#3);
}
\makeatother

% --- 文档开始 ---
\begin{document}

\begin{tikzpicture}
  % 1. 测试双点模式 (A,B)
  % 定义两个坐标
  \coordinate (P1) at (1,-1);
  \coordinate (P2) at (3,2);
  \segment[blue, dashed,thick](P1,P2)

  % % 2. 测试单点模式 (L)
  % 定义一个中心点
  \segment[orange,thick](P2)
\end{tikzpicture}

\end{document}